\chapter{แพลตฟอร์มการจัดการสวนทุเรียนอัจฉริยะและการเชื่อมต่อแดชบอร์ดข้อมูล}
\label{chap:ch07}

\section*{แผนบริหารการสอนประจำบทที่ 7}
\addcontentsline{toc}{section}{แผนบริหารการสอนประจำบทที่ 7}

\noindent\textbf{หัวข้อเนื้อหาประจำบท}
\begin{enumerate}
    \item สถาปัตยกรรมแดชบอร์ดและการรวมศูนย์ข้อมูลฟาร์ม
    \item การประมวลผลที่ขอบเครือข่ายและการตัดสินใจอัตโนมัติ
    \item การบูรณาการภาพถ่ายโดรนและดัชนีพืชพรรณ NDVI
\end{enumerate}

\noindent\textbf{วัตถุประสงค์เชิงพฤติกรรม}
\begin{enumerate}
    \item อธิบายหลักการ ทฤษฎี และสถาปัตยกรรมเทคโนโลยีประจำบทที่ 7 ได้อย่างถูกต้อง
    \item ประยุกต์ใช้เครื่องมือ อุปกรณ์เซนเซอร์ และระบบปัญญาประดิษฐ์เพื่อการจัดการสวนทุเรียนได้อย่างมีประสิทธิภาพ
    \item วิเคราะห์ สรุปผลข้อมูล และแก้ไขปัญหาเชิงฟิสิกส์เกษตรและคุณภาพดินได้อย่างเป็นระบบ
\end{enumerate}

\noindent\textbf{การวัดและประเมินผล}
\begin{table}[H]
\centering\small
\begin{tabularx}{\textwidth}{p{0.55\textwidth} p{0.25\textwidth} c}
\toprule
\textbf{เกณฑ์การประเมินผล} & \textbf{เครื่องมือวัดผล} & \textbf{สัดส่วนคะแนน} \\
\midrule
1. ทักษะการปฏิบัติงานและการใช้อุปกรณ์ & แบบประเมินทักษะภาคสนาม & 40\% \\
2. รายงานข้อมูลและการวิเคราะห์ผล & การตรวจสมุดบันทึกและดาต้าชีท & 30\% \\
3. แบบฝึกหัดและคำถามท้ายบท & แบบทดสอบท้ายบทเรียน & 20\% \\
4. การมีส่วนร่วมและความรับผิดชอบ & การเข้าเรียนและการตรงต่อเวลา & 10\% \\
\midrule
\textbf{รวมทั้งสิ้น} & & \textbf{100\%} \\
\bottomrule
\end{tabularx}
\end{table}
\newpage

% ==========================================================
% เนื้อหาบทเรียนประจำบทที่ 7
% ==========================================================

\section{สถาปัตยกรรมแดชบอร์ดและการรวมศูนย์ข้อมูลฟาร์ม}
\index{สถาปัตยกรรมแดชบอร์ดและการรวมศูนย์ข้อมูลฟาร์ม}

การศึกษาและทำความเข้าใจเกี่ยวกับสถาปัตยกรรมแดชบอร์ดและการรวมศูนย์ข้อมูลฟาร์ม (Centralized Farm Management Dashboard Architecture) มีความสำคัญอย่างยิ่งในระบบเกษตรอัจฉริยะ โดยมีรายละเอียดสาระสำคัญดังนี้

แพลตฟอร์มการจัดการสวนทุเรียนอัจฉริยะทำหน้าที่เป็นศูนย์กลางรวบรวม วิเคราะห์ และนำเสนอข้อมูลเชิงพื้นที่ (Spatial Data) และข้อมูลอนุกรมเวลา (Time-series Data) จากเซนเซอร์ภาคสนามทั้งหมด

องค์ประกอบของแดชบอร์ดอัจฉริยะ:
1. เครื่องมือแสดงผลภาพ (Visualization Widgets): แผนที่แปลงทุเรียนแบ่งโซน (GIS Farm Map), กราฟแนวโน้มความชื้นดินหลายระดับความลึก, มาตรวัดอุณหภูมิและความชื้นอากาศแบบเข็มดิจิทัล, และสถานะการทำงานของปั๊มน้ำ
2. ฐานข้อมูลอนุกรมเวลา (Time-Series Database): นิยมใช้ InfluxDB หรือ TimescaleDB ซึ่งได้รับการปรับแต่งมาเป็นพิเศษสำหรับการบันทึกและสืบค้นข้อมูลจากอุปกรณ์ IoT นับล้านจุดได้อย่างรวดเร็ว
3. แพลตฟอร์มแดชบอร์ดโอเพนซอร์ส: เช่น ThingsBoard, Grafana หรือ Node-RED Dashboard ซึ่งรองรับการปรับแต่งหน้าจอให้เหมาะสมกับการดูผ่านสมาร์ทโฟนและแท็บเล็ต

\begin{techbox}[ข้อกำหนดทางเทคนิคสถาปัตยกรรม AIoT]
การเชื่อมต่อระบบเซนเซอร์ไร้สายต้องรักษาความมั่นคงปลอดภัยของข้อมูล มีการเข้ารหัสระดับ AES-128 และบริหารจัดการพลังงานให้โหนดเซนเซอร์ทำงานได้อย่างต่อเนื่องไม่น้อยกว่า 3 ปี
\end{techbox}

\section{การประมวลผลที่ขอบเครือข่ายและการตัดสินใจอัตโนมัติ}
\index{การประมวลผลที่ขอบเครือข่ายและการตัดสินใจอัตโนมัติ}

การศึกษาและทำความเข้าใจเกี่ยวกับการประมวลผลที่ขอบเครือข่ายและการตัดสินใจอัตโนมัติ (Edge Computing and Automated Decision Making) มีความสำคัญอย่างยิ่งในระบบเกษตรอัจฉริยะ โดยมีรายละเอียดสาระสำคัญดังนี้

ในพื้นที่ห่างไกลที่สัญญาณอินเทอร์เน็ตอาจขัดข้องเป็นครั้งคราว สถาปัตยกรรม Edge Computing มีบทบาทสำคัญอย่างยิ่ง:
- คอนโทรลเลอร์เกตเวย์ในสวนจะรันอัลกอริทึมการตัดสินใจเฉพาะที่ (Local Rule Engine) สามารถสั่งเปิด-ปิดวาล์วน้ำได้เองโดยไม่ต้องพึ่งพาคลาวด์ตลอดเวลา
- เมื่ออินเทอร์เน็ตกลับมาเชื่อมต่อ ระบบจะทำการซิงโครไนซ์ข้อมูลย้อนหลัง (Store-and-Forward) ขึ้นสู่คลาวด์โดยอัตโนมัติ

\section{การบูรณาการภาพถ่ายโดรนและดัชนีพืชพรรณ NDVI}
\index{การบูรณาการภาพถ่ายโดรนและดัชนีพืชพรรณ NDVI}

การศึกษาและทำความเข้าใจเกี่ยวกับการบูรณาการภาพถ่ายโดรนและดัชนีพืชพรรณ NDVI (Multispectral Drone Imagery and NDVI Index) มีความสำคัญอย่างยิ่งในระบบเกษตรอัจฉริยะ โดยมีรายละเอียดสาระสำคัญดังนี้

นอกจากการตรวจวัดระดับจุดด้วยเซนเซอร์ในดินแล้ว การใช้โดรนติดกล้องหลายช่วงคลื่น (Multispectral Camera) ช่วยให้มองเห็นสุขภาพของต้นทุเรียนทั้งสวนในมุมกว้าง

ดัชนีพืชพรรณผลต่างปกติ (Normalized Difference Vegetation Index; NDVI) คำนวณจาก:
\[ \text{NDVI} = \frac{\text{NIR} - \text{Red}}{\text{NIR} + \text{Red}} \]
โดยที่ NIR คือการสะท้อนของคลื่นแสงใกล้อินฟราเรด (Near-Infrared: 840 nm) และ Red คือการสะท้อนของแสงสีแดง (660 nm)
- ต้นทุเรียนสมบูรณ์แข็งแรง: คลอโรฟิลล์ดูดกลืนแสงสีแดงสูง และโครงสร้างเซลล์ใบสะท้อน NIR สูง ทำให้ค่า NDVI อยู่ในช่วง 0.7 ถึง 0.9
- ต้นทุเรียนเริ่มแสดงอาการป่วย ขาดน้ำ หรือถูกโรครากเน่าคุกคาม: ค่า NDVI จะลดต่ำลงเหลือ 0.3 ถึง 0.5 ก่อนที่ตามนุษย์จะมองเห็น ช่วยให้เกษตรกรเข้าแก้ไขปัญหาได้อย่างทันท่วงที

\section{ขั้นตอนและวิธีปฏิบัติการทดลอง}
\index{ขั้นตอนการทดลองประจำบทที่ 7}

\subsection{การสร้างและกำหนดค่าหน้าจอแสดงผลบน ThingsBoard IoT Platform}
\begin{sensorbox}[การสร้างและกำหนดค่าหน้าจอแสดงผลบน ThingsBoard IoT Platform]
1. สมัครใช้งาน ThingsBoard Cloud หรือติดตั้งบนเซิร์ฟเวอร์ส่วนตัว
2. สร้าง Device Profile สำหรับชุดโหนดวัดสภาพแวดล้อมและโพรบวัดดิน
3. ตั้งค่า Rule Chain เพื่อประมวลผลข้อมูลขาเข้าและแปลงหน่วยอุณหภูมิเป็นองศาเซลเซียส
4. ออกแบบหน้า Dashboard: เพิ่ม Time-series Chart แสดงความชื้นดิน 3 ระดับ, เพิ่ม Card แสดงสถานะปั๊มน้ำ
5. สร้างปุ่มสวิตช์สั่งงานแบบ Manual Override สำหรับเปิดน้ำฉุกเฉิน
\end{sensorbox}

\subsection{การบินโดรนสำรวจแปลงทุเรียนและประมวลผลแผนที่ NDVI}
\begin{sensorbox}[การบินโดรนสำรวจแปลงทุเรียนและประมวลผลแผนที่ NDVI]
1. วางแผนเส้นทางการบินอัตโนมัติ (Mission Planning) ด้วยแอปพลิเคชัน กำหนดความสูง 50 เมตรเหนือพื้นดิน
2. ตั้งค่าระยะซ้อนทับของภาพ: ซ้อนทับด้านหน้า 75\% และซ้อนทับด้านข้าง 70\%
3. ทำการบินเก็บภาพถ่ายหลายช่วงคลื่นในช่วงเวลา 10:00 - 14:00 น. เพื่อหลีกเลี่ยงเงายาว
4. นำภาพเข้าโปรแกรมประมวลผลภาพถ่ายทางอากาศ (Photogrammetry) สร้างภาพโมเสกออร์โธ (Orthomosaic)
5. สร้างแผนที่สีดัชนี NDVI และระบุพิกัดต้นทุเรียนที่มีความสมบูรณ์ต่ำกว่าเกณฑ์
\end{sensorbox}

\subsection{การตั้งค่าระบบแจ้งเตือนผ่าน Telegram Bot อัตโนมัติ}
\begin{sensorbox}[การตั้งค่าระบบแจ้งเตือนผ่าน Telegram Bot อัตโนมัติ]
1. สนทนากับ BotFather บน Telegram เพื่อสร้างบอตใหม่และรับ API Token
2. ดึงค่า Chat ID ของกลุ่มผู้จัดการสวน
3. เขียนฟังก์ชันใน Node-RED เพื่อตรวจจับเหตุการณ์ผิดปกติ (เช่น ปั๊มน้ำไม่ทำงานตามเวลา)
4. ส่งข้อความแจ้งเตือนพร้อมแนบภาพกราฟและพิกัดต้นที่เกิดปัญหาเข้าสู่กลุ่มทันที
\end{sensorbox}

\section{ตารางบันทึกผลการทดลองและการอภิปรายผล}
\begin{table}[H]
\centering\small
\caption{ตารางบันทึกผลการตรวจวัดและข้อมูลพารามิเตอร์ประจำบทที่ 7}
\label{tab:ch07_datasheet}
\begin{tabularx}{\textwidth}{p{0.3\textwidth} p{0.3\textwidth} X}
\toprule
\textbf{พารามิเตอร์ / การทดสอบ} & \textbf{ค่าที่ตรวจวัดได้} & \textbf{การแปลผลและการวิเคราะห์} \\
\midrule
จุดตรวจวัดที่ 1 (บริเวณดินดอน) & ค่าความชื้นอยู่ในเกณฑ์ปกติ & ปริมาณน้ำเพียงพอต่อการเจริญ \\
จุดตรวจวัดที่ 2 (บริเวณที่ลุ่ม) & มีการสะสมความชื้นสูง & ควรชะลอการให้น้ำเพื่อป้องกันรากเน่า \\
สถานะสัญญาณโครงข่าย & RSSI: -85 dBm, SNR: +8 dB & คุณภาพสัญญาณ LoRaWAN ยอดเยี่ยม \\
\bottomrule
\end{tabularx}
\end{table}

\section{คำถามทบทวนและแบบฝึกหัดท้ายบท}
\begin{enumerate}
    \item จงอธิบายข้อดีของสถาปัตยกรรม Edge Computing เมื่อเปรียบเทียบกับระบบ Cloud-only ในการควบคุมระบบชลประทานสวนทุเรียน
    \item เหตุใดค่าดัชนี NDVI ของต้นทุเรียนจึงสามารถตรวจจับอาการผิดปกติของโรครากเน่าได้ล่วงหน้าก่อนที่มนุษย์จะสังเกตเห็นการเหี่ยวเฉาของใบ
    \item จงออกแบบโครงสร้างแดชบอร์ดสำหรับผู้จัดการสวนทุเรียนขนาด 50 ไร่ โดยระบุวิดเจ็ตและข้อมูลสำคัญที่ต้องแสดงผล
\end{enumerate}
